\documentclass[11pt]{article}

\usepackage[utf8]{inputenc}
\usepackage[T1]{fontenc}
\usepackage[spanish]{babel}
\usepackage{amsmath,amssymb,amsthm}
\usepackage[margin=2.7cm]{geometry}
\usepackage{enumitem}
\usepackage{hyperref}
\hypersetup{colorlinks=true, linkcolor=blue, urlcolor=blue, citecolor=blue}

% --- Operadores ---
\DeclareMathOperator{\tr}{tr}
\DeclareMathOperator{\im}{im}
\DeclareMathOperator{\Gr}{Gr}
\DeclareMathOperator{\Hess}{Hess}
\DeclareMathOperator{\sen}{sen}

% --- Entornos tipo teorema (numeración compartida por sección) ---
\theoremstyle{plain}
\newtheorem{teorema}{Teorema}[section]
\newtheorem{definicion}[teorema]{Definición}
\newtheorem{proposicion}[teorema]{Proposición}
\newtheorem{corolario}[teorema]{Corolario}
\newtheorem{lema}[teorema]{Lema}
\theoremstyle{definition}
\newtheorem{notacion}[teorema]{Notación}
\theoremstyle{remark}
\newtheorem{observacion}[teorema]{Observación}

\setlist[itemize]{itemsep=2pt, topsep=4pt}
\setlist[enumerate]{itemsep=2pt, topsep=4pt}

\title{\textbf{Resumen de Curvas y Superficies}\\[2mm]\large Definiciones, teoremas y resultados para el estudio}
\author{}
\date{}

\begin{document}
\maketitle
\vspace{-1cm}
\tableofcontents

\newpage

\section{Curvas en el plano y en el espacio}

\subsection{Curva, traza y regularidad}

\begin{definicion}
Una \textbf{curva} en el espacio es una aplicación diferenciable (de clase $C^\infty$) $\alpha:I\to\mathbb{R}^3$, donde $I\subseteq\mathbb{R}$ es un intervalo abierto. Si $\alpha(t)=(x(t),y(t),z(t))$, se llama \textbf{traza} de $\alpha$ al conjunto
\[
\tr\alpha=\im\alpha=\alpha(I).
\]
\end{definicion}

\begin{definicion}
Una curva $\alpha:I\to\mathbb{R}^3$ se dice \textbf{plana} si existe un plano $P\subset\mathbb{R}^3$ con $\tr\alpha\subset P$. Salvo un movimiento rígido se puede tomar $P=\{z=0\}$, de modo que una curva plana se identifica con una aplicación diferenciable $\alpha:I\to\mathbb{R}^2$.
\end{definicion}

\begin{observacion}
Las curvas no se exigen inyectivas: pueden tener autointersecciones. Una curva diferenciable puede tener ``picos'' en su traza (puntos donde la velocidad se anula), aunque la propia $\alpha$ sea diferenciable.
\end{observacion}

\begin{definicion}
Sea $\alpha:I\to\mathbb{R}^3$ una curva y $t\in I$ con $\alpha'(t)\neq0$. La \textbf{recta tangente} a $\alpha$ en $t$ es $R_t=\alpha(t)+\langle\alpha'(t)\rangle$. El punto $\alpha(t)$ se llama \textbf{regular} si $\alpha'(t)\neq0$ y \textbf{singular} si $\alpha'(t)=0$. La curva $\alpha$ es \textbf{regular} si todos sus puntos lo son, i.e. $\alpha'(t)\neq0\ \forall t\in I$.
\end{definicion}

\subsection{Cálculo vectorial útil}

Estas reglas de derivación y propiedades del producto escalar, vectorial y del determinante se usan constantemente al trabajar con curvas y superficies.

\begin{itemize}
\item Para $\alpha,\beta,\gamma:I\to\mathbb{R}^3$ derivables se cumple ``la regla del producto'' en cada caso:
\[
\langle\alpha,\beta\rangle'=\langle\alpha',\beta\rangle+\langle\alpha,\beta'\rangle, \qquad
(\alpha\land\beta)'=\alpha'\land\beta+\alpha\land\beta',
\]
\[
\frac{d}{dt}\det(\alpha,\beta,\gamma)=\det(\alpha',\beta,\gamma)+\det(\alpha,\beta',\gamma)+\det(\alpha,\beta,\gamma').
\]
\item $(|\alpha|^2)'=2\langle\alpha,\alpha'\rangle$, válida siempre; en cambio $|\alpha|'=\dfrac{\langle\alpha,\alpha'\rangle}{|\alpha|}$ solo si $\alpha\neq0$.
\item Interpretación geométrica: $\langle u,v\rangle=|u||v|\cos\theta$ (si $|v|=1$, es la proyección de $u$ sobre $v$); $u\land v$ es el único vector perpendicular a $u$ y $v$ tal que $\{u,v,u\land v\}$ está positivamente orientada y $|u\land v|=|u||v|\sen\theta$ (área del paralelogramo); $\det(w,u,v)=\langle w,u\land v\rangle$ es el volumen orientado del paralelepípedo. En coordenadas:
\[
u\land v=(u_2v_3-u_3v_2,\ u_3v_1-u_1v_3,\ u_1v_2-u_2v_1) = 
\begin{vmatrix}
    i & j & k\\
    u_1 & u_2 & u_3 \\
    v_1 & v_2 & v_3
\end{vmatrix}.
\]
\item Propiedades útiles, para $u,v,w\in\mathbb{R}^3$:
\begin{itemize}
\item Si $A\in O(3)$: $\langle Au,Av\rangle=\langle u,v\rangle$ y $Au\land Av=\det(A)\,A(u\land v)$.
\item Si $A\in M_3(\mathbb{R})$: $\det(Au,Av,Aw)=\det(A)\det(u,v,w)$.
\item $\langle u\land v,w\rangle=\det(u,v,w)$.
\item $\land$ es bilineal y antisimétrico ($u\land v=-(v\land u)$); si $u,v$ son linealmente dependientes, $u\land v=0$.
\item En $\mathbb{R}^2$: $|\det(u,v)|=|u||v|\sen\theta$.
\end{itemize}
\end{itemize}

\subsection{Longitud de una curva}

\begin{definicion}
Sea $\alpha:I\to\mathbb{R}^3$ una curva y $[a,b]\subset I$. La \textbf{longitud} de $\alpha$ entre $a$ y $b$ es
\[
L_a^b(\alpha)=\int_a^b|\alpha'(t)|\,dt.
\]
\end{definicion}

\begin{observacion}
Esta cantidad es el límite (y el supremo) de las longitudes de las poligonales $L_a^b(\alpha,P)=\sum_i|\alpha(t_i)-\alpha(t_{i-1})|$ asociadas a particiones $P$ de $[a,b]$, cuando el diámetro de $P$ tiende a $0$. Además siempre se tiene $|\alpha(b)-\alpha(a)|\le L_a^b(\alpha,P)\le L_a^b(\alpha)$.
\end{observacion}

\subsection{Reparametrizaciones, deformaciones y movimientos rígidos}

\begin{observacion}
Si $h:J\to I$ ($J$ intervalo abierto) es diferenciable y mantiene la regularidad de una curva regular $\alpha$ (es decir, $\beta=\alpha\circ h$ es regular), entonces $h$ ha de ser un \textbf{difeomorfismo} (creciente o decreciente).
\end{observacion}

\begin{definicion}[Reparametrización]
Sea $\alpha:I\to\mathbb{R}^3$ una curva y $h:J\to I$ un difeomorfismo sobreyectivo. La curva $\beta=\alpha\circ h$ se llama \textbf{reparametrización} de $\alpha$ mediante $h$; se tiene $\tr\beta=\tr\alpha$. Si $\alpha$ es regular y $h'\neq0$, entonces $\beta$ también es regular.
\end{definicion}

\begin{proposicion}
La longitud es invariante por reparametrizaciones: si $\beta=\alpha\circ h$ y $[a,b]\subset J$,
\[
L_a^b(\beta)=L_{h(a)}^{h(b)}(\alpha).
\]
\end{proposicion}

\begin{observacion}[Deformaciones]
Si $\Phi:\mathbb{R}^3\to\mathbb{R}^3$ es diferenciable y $\beta=\Phi\circ\alpha$, entonces $\beta'(t)=(d\Phi)_{\alpha(t)}(\alpha'(t))$. Para que $\beta$ sea regular siempre que $\alpha$ lo sea, basta con que $\Phi$ sea un \textbf{difeomorfismo local}. En general, ni la traza ni la longitud se conservan bajo deformaciones.
\end{observacion}

\begin{definicion}[Movimiento rígido]
Una aplicación $M:\mathbb{R}^3\to\mathbb{R}^3$ dada por $M(x)=Ax+b$, con $A\in O(3)$ y $b\in\mathbb{R}^3$, se llama \textbf{movimiento rígido}; es directo si $\det A=1$ e inverso si $\det A=-1$. Conserva ángulos y distancias.
\end{definicion}

\begin{proposicion}
Los movimientos rígidos conservan la longitud de las curvas: $L_a^b(M\circ\alpha)=L_a^b(\alpha)$.
\end{proposicion}

\subsection{Curvas parametrizadas por el arco}

\begin{definicion}
Una curva $\alpha:I\to\mathbb{R}^3$ está \textbf{parametrizada por el arco} (p.p.a.) si $|\alpha'(t)|=1\ \forall t\in I$.
\end{definicion}

\begin{proposicion}
Toda curva regular admite una reparametrización p.p.a. (tomando $h={(S_a)}^{-1}$, donde $S_a(t)=\int_a^t|\alpha'(u)|\,du$ es la \emph{longitud orientada} respecto de $a\in I$).
\end{proposicion}

\begin{proposicion}
Las reparametrizaciones p.p.a. de una curva p.p.a. $\alpha:I\to\mathbb{R}^3$ son exactamente las de la forma $h(t)=\pm t+c$, con $c\in\mathbb{R}$.
\end{proposicion}

\begin{observacion}[Interés geométrico]
Si $\alpha$ está p.p.a., entonces $|\alpha'(s)|^2=1\ \forall s$ y, derivando, $\alpha''(s)\perp\alpha'(s)\ \forall s\in I$: la aceleración de una curva de celeridad constante únicamente puede ``rotar'' la velocidad, nunca cambiar su módulo. Esto justifica restringirse a curvas p.p.a. para el estudio puramente geométrico (no físico) de las curvas.
\end{observacion}

\subsection{Teoría local de curvas planas: diedro de Frenet y curvatura}

Sea $\alpha:I\to\mathbb{R}^2$ una curva plana, regular y p.p.a.; escribimos $T(s)=\alpha'(s)$.

\begin{definicion}
El \textbf{vector normal} de $\alpha$ en $s$ es $N(s)=JT(s)$, donde $J(x,y)=(-y,x)$ es el giro de $90^\circ$ en sentido antihorario. La base ortonormal positivamente orientada $\{T(s),N(s)\}$ se llama \textbf{diedro de Frenet} de $\alpha$ en $s$.
\end{definicion}

\begin{proposicion}[Ecuaciones de Frenet, caso plano]
Existe una única función diferenciable $k:I\to\mathbb{R}$, la \textbf{curvatura} de $\alpha$, tal que
\[
\boxed{T'(s)=k(s)N(s), \qquad N'(s)=-k(s)T(s)}.
\]
Además $k(s)=\langle T'(s),N(s)\rangle$ y $|k(s)|=|\alpha''(s)|$.
\end{proposicion}

\begin{observacion}[Interpretación]
El signo de $k$ indica el sentido de giro de la curva: $k(s)>0$ si la curva gira a la izquierda (la aceleración $\alpha''(s)$ tiene el mismo sentido que $N(s)$), $k(s)<0$ si gira a la derecha, y $k(s)=0$ en los puntos sin curvatura (la aceleración es nula).
\end{observacion}

\begin{proposicion}[Fórmula general de la curvatura]
Si $\alpha:I\to\mathbb{R}^2$ es regular (no necesariamente p.p.a.), su curvatura viene dada por
\[
\boxed{k_\alpha(t)=\frac{\langle J\alpha'(t),\alpha''(t)\rangle}{|\alpha'(t)|^3}=\frac{\det(\alpha'(t),\alpha''(t))}{|\alpha'(t)|^3}}.
\]
Esta fórmula no depende de la reparametrización por el arco utilizada para definir $k_\alpha$, y no puede extenderse a puntos singulares.
\end{proposicion}

\begin{proposicion}[Curvatura, reparametrizaciones y movimientos rígidos]\
\begin{itemize}
\item Si $\beta(t)=\alpha(h(t))$ con $h$ creciente, $k_\beta=k_\alpha\circ h$; si $h$ es decreciente, $k_\beta=-k_\alpha\circ h$.
\item Si $M(x)=Ax+b$ es un movimiento rígido, $k_{M\circ\alpha}=\det(A)\,k_\alpha$: la curvatura se conserva si $M$ es directo y cambia de signo si $M$ es inverso.
\end{itemize}
\end{proposicion}

\begin{observacion}[Resultados de referencia]\
\begin{itemize}
\item Las rectas (segmentos de recta) tienen curvatura idénticamente nula.
\item Una circunferencia de radio $r$ tiene curvatura constante $k\equiv\tfrac1r$ (decrece al aumentar el radio).
\item Las únicas curvas p.p.a. de curvatura constante son la recta ($k\equiv0$) y la circunferencia de radio $1/|k_0|$ ($k\equiv k_0\neq0$).
\item Para la curva asociada a la gráfica de $f:I\to\mathbb{R}$ derivable, $\alpha_f(t)=(t,f(t))$, se tiene
\[
k_{\alpha_f}(t)=\frac{f''(t)}{(1+f'(t)^2)^{3/2}}:
\]
así, concavidad/convexidad de $f$ corresponden a curvatura negativa/positiva, y los puntos de inflexión a curvatura nula.
\end{itemize}
\end{observacion}

\begin{teorema}[Fundamental de la Teoría Local de Curvas Planas]
Sea $I\subset\mathbb{R}$ un intervalo abierto y $k_0:I\to\mathbb{R}$ una función diferenciable. Entonces existe una curva p.p.a. $\alpha:I\to\mathbb{R}^2$ tal que $k_\alpha(s)=k_0(s)\ \forall s\in I$. Además, $\alpha$ es única salvo movimientos rígidos directos.
\end{teorema}

\begin{observacion}
Este teorema dice que, salvo su posición en el plano, una curva plana queda determinada por su curvatura en función del parámetro arco. En particular, si $\alpha,\beta:I\to\mathbb{R}^2$ son curvas p.p.a. con $k_\alpha=-k_\beta$, existe un movimiento rígido \emph{inverso} $M$ con $\beta=M\circ\alpha$.
\end{observacion}

\subsection{Teoría local de curvas en el espacio: triedro de Frenet y torsión}

Sea $\alpha:I\to\mathbb{R}^3$ p.p.a.; escribimos $T(s)=\alpha'(s)$.

\begin{definicion}
La \textbf{curvatura} de $\alpha$ en $s$ es $k(s)=|T'(s)|=|\alpha''(s)|\ge0$.
\end{definicion}

A partir de aquí se supone $k(s)>0\ \forall s\in I$ (una ``regularidad de orden $2$''), lo que permite definir el resto del triedro.

\begin{definicion}
El \textbf{vector normal} de $\alpha$ en $s$ es $N(s)=\dfrac{T'(s)}{k(s)}$, y el \textbf{vector binormal} es $B(s)=T(s)\land N(s)$. La base ortonormal positivamente orientada $\{T(s),N(s),B(s)\}$ es el \textbf{triedro de Frenet} de $\alpha$ en $s$.
\end{definicion}

\begin{definicion}
La \textbf{torsión} de $\alpha$ en $s$ es el único escalar $\tau(s)$ tal que $B'(s)=\tau(s)N(s)$.
\end{definicion}

\begin{proposicion}[Ecuaciones de Frenet, caso espacial]
\[
\boxed{T'=kN, \qquad N'=-kT-\tau B, \qquad B'=\tau N}.
\]
\end{proposicion}

\begin{observacion}
La torsión mide la tendencia de la curva a salirse del \emph{plano osculador} (el generado por $T$ y $N$ en cada instante); su signo indica el sentido en que se produce esa salida.
\end{observacion}

\begin{proposicion}
Si $\alpha$ es plana, su curvatura como curva del espacio coincide en valor absoluto con su curvatura como curva plana: $k_\alpha^{\mathbb{R}^3}=|k_\alpha^{P}|$. En particular:
\begin{itemize}
\item Una recta no tiene vector normal definido ($k\equiv0$).
\item Una circunferencia de radio $r$ tiene $k\equiv\tfrac1r$ y $\tau\equiv0$.
\end{itemize}
\end{proposicion}

\begin{teorema}
Sea $\alpha:I\to\mathbb{R}^3$ p.p.a. con $k(s)>0\ \forall s$. Entonces $\alpha$ es una curva plana si y solo si $\tau\equiv0$.
\end{teorema}

\begin{observacion}[Hélice circular recta]
Para $a,b\in\mathbb{R}$ con $ab\neq0$, la curva
\[
\alpha(s)=\left(a\cos\frac{s}{\sqrt{a^2+b^2}},\ a\sen\frac{s}{\sqrt{a^2+b^2}},\ \frac{bs}{\sqrt{a^2+b^2}}\right)
\]
se llama \textbf{hélice circular recta}; está p.p.a., no es plana, y tiene curvatura y torsión \emph{constantes}:
\[
\boxed{k=\frac{|a|}{a^2+b^2}, \qquad \tau=\frac{-b}{a^2+b^2}}.
\]
El valor $|a|$ es el radio del cilindro que contiene a la hélice, y el signo de $\tau$ determina el sentido de giro de la misma (levógira o dextrógira).
\end{observacion}

\begin{teorema}[Caracterización de la hélice]
Sea $\alpha:I\to\mathbb{R}^3$ p.p.a. con curvatura $k>0$ constante y torsión $\tau$ constante. Entonces $\tr\alpha$ coincide, salvo un movimiento rígido, con la traza de una hélice circular recta (la recta y la circunferencia son los casos límite $b=0$ ó $a=0$, respectivamente).
\end{teorema}

\begin{proposicion}[Curvas en una esfera]
Sea $\alpha:I\to\mathbb{R}^3$ p.p.a. con $k>0$ y $\tau>0$, con $\tr\alpha$ contenida en una esfera $S^2(a,R)$. Entonces
\[
R^2=\frac{1}{k^2}+\frac{1}{\tau^2}\left[\left(\frac1k\right)'\right]^2.
\]
Además, si $\tau\equiv\tau_0$ es constante, existen $b,c\in\mathbb{R}$ tales que
\[
k(s)=\frac{1}{b\cos(|\tau_0|s)+c\sen(|\tau_0|s)}.
\]
\end{proposicion}

\begin{proposicion}
Sea $\alpha:I\to\mathbb{R}^3$ p.p.a. con $k>0$. Entonces $\alpha$ es un arco de circunferencia si y solo si $k$ es constante y $\tr\alpha$ está contenida en una esfera.
\end{proposicion}

\begin{proposicion}[Curvatura, torsión y movimientos rígidos]
Si $M$ es un movimiento rígido y $\beta=M\circ\alpha$ (con $\alpha$ p.p.a.), entonces
\[
k_\beta=k_\alpha, \qquad \tau_\beta=\begin{cases}\tau_\alpha & \text{si $M$ es directo}\\ -\tau_\alpha & \text{si $M$ es inverso.}\end{cases}
\]
\end{proposicion}

\begin{proposicion}[Fórmulas generales para curvas no p.p.a.]
Si $\alpha:I\to\mathbb{R}^3$ es regular (no necesariamente p.p.a.) y $k_\alpha>0$:
\[
\boxed{k_\alpha(t)=\frac{|\alpha'(t)\land\alpha''(t)|}{|\alpha'(t)|^3}}, \qquad
\boxed{\tau_\alpha(t)=-\frac{\det(\alpha'(t),\alpha''(t),\alpha'''(t))}{|\alpha'(t)\land\alpha''(t)|^2}}.
\]
\end{proposicion}

\begin{teorema}[Fundamental de la Teoría Local de Curvas en el Espacio]
Sea $I\subset\mathbb{R}$ un intervalo abierto y $k_0,\tau_0:I\to\mathbb{R}$ funciones diferenciables con $k_0>0$. Entonces existe una curva p.p.a. $\alpha:I\to\mathbb{R}^3$ tal que $k_\alpha=k_0$ y $\tau_\alpha=\tau_0$ en $I$. Además, $\alpha$ es única salvo un movimiento rígido directo.
\end{teorema}


\newpage
\section{Superficies en el espacio}

\subsection{Definición de superficie}

\begin{definicion}
Una \textbf{superficie} en $\mathbb{R}^3$ es un subconjunto $S\subset\mathbb{R}^3$ no vacío tal que para todo $p\in S$ existen un abierto $U\subset\mathbb{R}^2$, un entorno abierto $V$ de $p$ en $S$ y una aplicación diferenciable $X:U\to\mathbb{R}^3$ de modo que:
\begin{enumerate}
\item $X(U)=V$;
\item $X:U\to V$ es un homeomorfismo;
\item $(dX)_q:\mathbb{R}^2\to\mathbb{R}^3$ es inyectiva para todo $q\in U$.
\end{enumerate}
\end{definicion}

\begin{notacion}
A $X$ se le llama \textbf{parametrización} (local), \textbf{mapa} o \textbf{carta}; a $(u,v)\in U$, \textbf{coordenadas locales}; a $V$, \textbf{entorno parametrizado}; y a $U$, \textbf{dominio de la parametrización}. Se escribe $X_u=\partial X/\partial u$, $X_v=\partial X/\partial v$.
\end{notacion}

\begin{observacion}
La condición 3 equivale a que $X_u(q)$ y $X_v(q)$ sean linealmente independientes para todo $q\in U$, o equivalentemente, a que $X_u(q)\land X_v(q)\neq0\ \forall q\in U$.
\end{observacion}

\subsection{Construcción de superficies}

\begin{teorema}[de la función implícita]
Sea $O\subset\mathbb{R}^3$ abierto, $f:O\to\mathbb{R}$ diferenciable, $a\in\mathbb{R}$ y $p=(x_0,y_0,z_0)\in f^{-1}(\{a\})$ con $f_z(p)\neq0$. Entonces existen $U\subset\mathbb{R}^2$ entorno de $(x_0,y_0)$, $V\subset\mathbb{R}$ entorno de $z_0$ y $g:U\to V$ diferenciable, con $g(x_0,y_0)=z_0$ y
\[
\{p\in U\times V: f(p)=a\}=\{(x,y,g(x,y)):(x,y)\in U\}.
\]
\end{teorema}

\begin{definicion}
Para $f:O\subset\mathbb{R}^3\to\mathbb{R}$ diferenciable, $a\in\mathbb{R}$ es \textbf{valor regular} de $f$ si $(df)_p\neq0$ (equivalentemente, $\nabla f(p)\neq0$) para todo $p\in f^{-1}(\{a\})$.
\end{definicion}

\begin{proposicion}\label{prop:valor_regular}
Si $f:O\subset\mathbb{R}^3\to\mathbb{R}$ es diferenciable y $a\in\mathbb{R}$ es un valor regular de $f$, entonces $S=f^{-1}(\{a\})$ es vacío o una superficie.
\end{proposicion}

\noindent Esta proposición es la principal herramienta para construir superficies a partir de ecuaciones implícitas.

\subsection{Catálogo de superficies y propiedades de estabilidad}

\begin{itemize}
\item Los \textbf{planos} y, en general, la \textbf{gráfica} de cualquier función diferenciable $f:O\subset\mathbb{R}^2\to\mathbb{R}$ son superficies, cubiertas por una única parametrización $X(u,v)=(u,v,f(u,v))$.
\item Si $S$ es una superficie y $\emptyset\neq O\subset S$ es abierto (relativo a $S$), entonces $O$ es una superficie. En consecuencia, toda componente conexa de una superficie es una superficie.
\item Si $S=\bigcup_{i\in I} S_i$ con cada $S_i$ abierto en $S$ y superficie, entonces $S$ es una superficie. En particular, una unión \underline{disjunta} de superficies es una superficie.
\item Si $S$ es una superficie y $\phi:O\to O'$ es un difeomorfismo entre abiertos de $\mathbb{R}^3$ con $S\subset O$, entonces $\phi(S)$ es una superficie.
\item Las \textbf{esferas} $S^2(c,R)$ y los \textbf{elipsoides} son superficies.
\item Las \textbf{cuádricas} sin puntos singulares (elipsoides, hiperboloides, paraboloides, cilindros) son superficies, por la Proposición \ref{prop:valor_regular}. Los conos y los pares de planos secantes quedan excluidos.
\item El \textbf{toro de revolución} es una superficie.
\item El \textbf{cono} $\{x^2+y^2=z^2\}$ \emph{no} es una superficie: cualquier entorno del vértice menos el vértice tiene dos componentes conexas, mientras que un abierto de $\mathbb{R}^2$ menos un punto sigue siendo conexo, así que no puede haber homeomorfismo local en ese punto.
\item Una superficie \textbf{compacta} tiene un número finito de componentes conexas.
\end{itemize}

\subsection{Aplicaciones diferenciables}

\begin{definicion}
Sean $S,S'\subset\mathbb{R}^3$ superficies y $O\subset\mathbb{R}^n$ abierto.
\begin{enumerate}[label=\alph*)]
\item $f:S\to\mathbb{R}^m$ es diferenciable si $f\circ X:U\to\mathbb{R}^m$ es diferenciable (en el sentido del análisis) para \emph{toda} parametrización $X:U\to S$.
\item $f:O\to S$ es diferenciable si $i\circ f:O\to\mathbb{R}^3$ lo es, donde $i:S\hookrightarrow\mathbb{R}^3$.
\item $f:S\to S'$ es diferenciable si $i'\circ f:S\to\mathbb{R}^3$ lo es en el sentido a), donde $i':S'\hookrightarrow\mathbb{R}^3$.
\end{enumerate}
\end{definicion}

\begin{proposicion}
Toda aplicación diferenciable (en cualquiera de los tres sentidos anteriores) es continua.
\end{proposicion}

\begin{proposicion}
Si $f,g:S\to\mathbb{R}^m$ y $\lambda:S\to\mathbb{R}$ son diferenciables, también lo son $f+g$ y $\lambda f$. Además $\langle f,g\rangle:S\to\mathbb{R}$ es diferenciable y, si $m=3$, también lo es $f\land g:S\to\mathbb{R}^3$.
\end{proposicion}

\begin{lema}[Cambio de cartas]\label{lema:cambio_cartas}
Si $X_1:U_1\to S$ y $X_2:U_2\to S$ son parametrizaciones con $X_1(U_1)\cap X_2(U_2)\neq\emptyset$, entonces el \textbf{cambio de coordenadas} $X_2^{-1}\circ X_1$ (definido entre los correspondientes abiertos de $\mathbb{R}^2$) es un \textbf{difeomorfismo}.
\end{lema}

\begin{corolario}\label{coro:recubrimiento}
Sea $f:S\to\mathbb{R}^m$. Si para todo $p\in S$ existe una parametrización $X_p:U_p\to S$ con $p\in X_p(U_p)$ y tal que $f\circ X_p$ es diferenciable, entonces $f$ es diferenciable. Es decir: \underline{no es necesario comprobar la definición para todas las parametrizaciones, basta un recubrimiento de $S$}.
\end{corolario}

\begin{observacion}[Propiedades adicionales]\
\begin{itemize}
\item Toda parametrización $X:U\to S$ es diferenciable, y $X:U\to X(U)$ es un difeomorfismo (ver Sección \ref{sec:difeo}).
\item \textbf{Restricciones}: si $S\subset O$ con $O\subset\mathbb{R}^3$ abierto y $F:O\to\mathbb{R}^m$ es diferenciable, $f=F|_S$ es diferenciable. Esto proporciona muchos ejemplos (funciones constantes, la inclusión $i:S\hookrightarrow\mathbb{R}^3$, la distancia al cuadrado a un punto fijo, etc.).
\item \textbf{Regla de la cadena}: si $f_1:S_1\to S_2$ y $f_2:S_2\to S_3$ son diferenciables (con $S_i$ superficies o abiertos euclídeos), $f_2\circ f_1$ es diferenciable.
\item Si $S=\bigcup_{i\in I}S_i$ con cada $S_i\subset S$ abierto y $f|_{S_i}$ diferenciable para todo $i$, entonces $f$ es diferenciable.
\end{itemize}
\end{observacion}

\subsection{Difeomorfismos}\label{sec:difeo}

\begin{definicion}
Si $S,S'$ son superficies, $f:S\to S'$ es un \textbf{difeomorfismo} si es diferenciable, biyectiva y su inversa $f^{-1}$ también es diferenciable. Dos superficies son \textbf{difeomorfas} si existe un difeomorfismo entre ellas.
\end{definicion}

\begin{observacion}\
\begin{itemize}
\item Toda parametrización $X:U\to S$ es un difeomorfismo entre $U$ (abierto de $\mathbb{R}^2$) y $X(U)$.
\item ``Ser difeomorfas'' es una relación de equivalencia entre superficies.
\item Dos superficies difeomorfas son homeomorfas. El recíproco es cierto para superficies (un teorema profundo de Geometría Diferencial), aunque en general no lo es para variedades de dimensión mayor.
\end{itemize}
\end{observacion}

\subsection{Plano tangente}

\begin{definicion}
Sea $S$ una superficie y $p\in S$. Un vector $v\in\mathbb{R}^3$ es \textbf{tangente a $S$ en $p$} si existe una curva diferenciable $\alpha:\,]-\varepsilon,\varepsilon[\,\to S$ con $\alpha(0)=p$ y $\alpha'(0)=v$. El conjunto de todos ellos se denota $T_pS\subset\mathbb{R}^3$.
\end{definicion}

\begin{teorema}
Si $X:U\to S$ es una parametrización con $p\in X(U)$, entonces
\[
\boxed{T_pS=(dX)_{X^{-1}(p)}(\mathbb{R}^2)}.
\]
En particular, $T_pS$ es un \textbf{plano vectorial} (subespacio de dimensión $2$) que no depende de $X$, y $\{X_u,X_v\}$ (evaluados en $X^{-1}(p)$) es una base de $T_pS$. A $T_pS$ se le llama \textbf{plano tangente} a $S$ en $p$.
\end{teorema}

\begin{proposicion}[Caso implícito]
Si $a$ es valor regular de $f:O\subset\mathbb{R}^3\to\mathbb{R}$ y $S=f^{-1}(\{a\})$, entonces para $p\in S$:
\[
\boxed{T_pS=\ker(df)_p=\langle\nabla f(p)\rangle^\perp}.
\]
\end{proposicion}

\begin{definicion}
La \textbf{recta normal} a $S$ en $p$ es la recta afín que pasa por $p$ y es perpendicular a $T_pS$.
\end{definicion}

\begin{observacion}
Si $S'\subset S$ es un abierto relativo de una superficie $S$, y $p\in S'$, entonces $T_pS'=T_pS$: el plano tangente es un objeto local.
\end{observacion}

\subsection{Diferencial de una aplicación diferenciable}

\begin{definicion}
Sea $S$ una superficie y $f:S\to\mathbb{R}^m$ diferenciable (tipo a)). Para $p\in S$ y $v\in T_pS$, eligiendo $\alpha:\,]-\varepsilon,\varepsilon[\,\to S$ con $\alpha(0)=p,\ \alpha'(0)=v$, se define la \textbf{diferencial}
\[
(df)_p(v)=(f\circ\alpha)'(0).
\]
\end{definicion}

\begin{proposicion}
$(df)_p(v)$ no depende de la curva $\alpha$ elegida (solo de $v$), y $(df)_p:T_pS\to\mathbb{R}^m$ es una aplicación \textbf{lineal}.
\end{proposicion}

\begin{definicion}
Análogamente, si $f:O\to S$ (tipo b)) es diferenciable, su diferencial en $x\in O$ es $(df)_x:\mathbb{R}^n\to T_{f(x)}S$ (la diferencial usual, cuya imagen cae automáticamente en $T_{f(x)}S$). Si $f:S\to S'$ (tipo c)) es diferenciable, su diferencial en $p\in S$ es $(df)_p:T_pS\to T_{f(p)}S'$ (la diferencial de tipo a), cuya imagen cae automáticamente en $T_{f(p)}S'$).
\end{definicion}

\begin{observacion}[Diferenciales útiles]
Para $F:\mathbb{R}^3\to\mathbb{R}^m$ diferenciable, $S$ una superficie y $f=F|_S$, se tiene $(df)_p=(dF)_p|_{T_pS}$. En particular:
\begin{itemize}
\item $f(p)=|p-p_0|^2$: $\ (df)_p(v)=2\langle p-p_0,v\rangle$.
\item $f(p)=\langle p-p_0,u\rangle$ con $|u|=1$ (distancia orientada a un plano): $\ (df)_p(v)=\langle v,u\rangle$.
\item Inclusión $i:S\hookrightarrow\mathbb{R}^3$: $\ (di)_p$ es la inclusión de $T_pS$ en $\mathbb{R}^3$.
\item $f$ constante: $\ (df)_p\equiv0$.
\end{itemize}
\end{observacion}

\subsection{Propiedades importantes de la diferencial}

\begin{enumerate}
\item \textbf{Regla de la cadena}: si $f:S\to S'$ y $g:S'\to S''$ son diferenciables, $d(g\circ f)_p=(dg)_{f(p)}\circ(df)_p$ para todo $p\in S$.

\noindent\emph{Consecuencia}: si $\phi:S_1\to S_2$ es un difeomorfismo, $(d\phi)_p:T_pS_1\to T_{\phi(p)}S_2$ es un isomorfismo lineal y $(d\phi)_p^{-1}=(d\phi^{-1})_{\phi(p)}$.

\item \textbf{Extremos locales}: si $f:S\to\mathbb{R}$ tiene un extremo local en $p_0\in S$, entonces $(df)_{p_0}\equiv0$.

\item \textbf{Diferencial nula}: si $S$ es conexa, $f:S\to S'$ es diferenciable y $(df)_p\equiv0\ \forall p\in S$, entonces $f$ es constante.

\item \textbf{Teorema de la función inversa}: si $f:S\to S'$ es diferenciable y $(df)_{p_0}$ es un isomorfismo, existen entornos $V$ de $p_0$ en $S$ y $W$ de $f(p_0)$ en $S'$ con $f(V)=W$ y $f|_V:V\to W$ difeomorfismo.

\noindent\emph{Consecuencia}: si $(df)_p$ es isomorfismo para todo $p\in S$, $f$ es un \textbf{difeomorfismo local}; en particular $f$ es una aplicación abierta, luego $f(S)$ es abierto en $S'$.
\end{enumerate}

\begin{corolario}
Si $S'\subset S$ son dos superficies, entonces $S'$ es necesariamente un \textbf{abierto} de $S$ (no puede haber una superficie ``estrictamente menor'' contenida en otra sin ser abierta en ella).
\end{corolario}

\subsection{Otros resultados geométricos útiles}

\begin{itemize}
\item Toda superficie $S\subset\mathbb{R}^3$ tiene \textbf{interior vacío} en $\mathbb{R}^3$.
\item Para $S$ superficie y $p_0\in\mathbb{R}^3$, los puntos críticos de $f(p)=|p-p_0|^2$ sobre $S$ son exactamente los $p\in S$ tales que $p_0$ está en la recta normal a $S$ en $p$.
\item Si $S$ es compacta, por cada punto $p_0\in\mathbb{R}^3$ pasa al menos una recta normal a $S$.
\item Si $S$ es conexa y todas sus rectas normales son concurrentes en un punto $p_0$, entonces $S$ está contenida en una esfera centrada en $p_0$ (de hecho, es un abierto conexo de dicha esfera).
\item Si $S$ es cerrada, la distancia de $S$ a cualquier punto de $\mathbb{R}^3$ se alcanza (existe un punto que minimiza la distancia).
\item Si $S$ es compacta, existe una recta de $\mathbb{R}^3$ que corta a $S$ perpendicularmente en al menos dos de sus puntos.
\end{itemize}

\section{Segunda forma fundamental. Curvaturas}

\subsection{Orientabilidad}

\begin{definicion}
Diremos que la superficie $S$ es \textbf{orientable} si existe $N:S\to \mathbb{R}^3$ diferenciable tal que
\[
\forall p\in S \quad |N(p)| = 1 \quad \text{y}\quad N(p)\perp T_pS.
\]
Si $S$ es orientable, está \textbf{orientada} cuando se ha elegido tal $N$. A $N$ se le llama también \textbf{campo normal unitario} o \textbf{aplicación de Gauss}.
\end{definicion}

\begin{proposicion}
Sea $S$ una superficie conexa y orientable. Entonces existen exactamente dos orientaciones sobre $S$: si $N$ y $N'$ son dos orientaciones, necesariamente $N' = N$ o $N' = -N$.
\end{proposicion}

\begin{proposicion}
Toda superficie $S$ es localmente orientable: cada $p\in S$ admite un entorno $V$ en $S$ orientable. Si $X:U\to S$ es una parametrización con $p\in X(U)$, la aplicación de Gauss local es
\[
N^X = \frac{X_u\land X_v}{|X_u\land X_v|}.
\]
\end{proposicion}

\begin{observacion}[Ejemplos de superficies orientables]\
\begin{itemize}
\item Todo plano afín $P$ es orientable (el vector normal unitario constante orienta $P$).
\item Toda esfera $S^2(p_0,R)$ es orientable, con $N_p = \dfrac{p-p_0}{R}$.
\item Todo grafo $S = \Gr h$ de una función $h:U\subset\mathbb{R}^2\to\mathbb{R}$ es orientable (lo cubre una sola parametrización).
\item Las imágenes inversas de valores regulares de $f:O\subset\mathbb{R}^3\to\mathbb{R}$ son orientables, con $N = \dfrac{\nabla f}{|\nabla f|}$.
\item \textbf{Teorema de Brouwer-Samelson.} Toda superficie $S\subset\mathbb{R}^3$ compacta (o cerrada) es orientable.
\end{itemize}
\end{observacion}

\subsection{Segunda forma fundamental y curvaturas}

Sea $S$ una superficie orientada por $N:S\to S^2$. Para cada $p\in S$, la diferencial $(dN)_p:T_pS\to T_{N_p}S^2 = T_pS$ es un endomorfismo de $T_pS$. Sus invariantes geométricos son:

\begin{definicion}
La \textbf{curvatura de Gauss} y la \textbf{curvatura media} de $S$ en $p$ se definen como
\[
K(p) = \det(dN)_p, \qquad H(p) = -\frac{1}{2}\tr(dN)_p.
\]
El \textbf{endomorfismo de Weingarten} es $A_p = -(dN)_p$, con lo que
\[
K(p) = \det A_p, \qquad H(p) = \frac{1}{2}\tr A_p.
\]
La \textbf{segunda forma fundamental} de $S$ en $p$ (para la orientación $N$) es la forma bilineal $\sigma_p:T_pS\times T_pS\to\mathbb{R}$ dada por
\[
\sigma_p(v,w) = \langle A_p(v),w\rangle = -\langle (dN)_p(v),w\rangle.
\]
La \textbf{primera forma fundamental} $g_p$ de $S$ en $p$ es el producto escalar de $\mathbb{R}^3$ restringido a $T_pS$:
\[
g_p(v,w) = \langle v,w\rangle.
\]
\end{definicion}

\begin{observacion}
Al cambiar la orientación de $N$ a $-N$: la segunda forma fundamental y la curvatura media cambian de signo, pero la curvatura de Gauss $K$ no varía.
\end{observacion}

\begin{observacion}[Intuición geométrica de $\sigma_p$]
Para $v\in T_pS$ con $|v|=1$, la intersección del plano $P_v = p + \langle v, N_p\rangle$ con $S$ es la traza de una curva regular $\alpha$ con $\alpha(0)=p$, $\alpha'(0)=v$, llamada \textbf{sección normal} a $S$ en $p$. Se tiene
\[
\sigma_p(v,v) = \langle \alpha''(0), N_p\rangle = k_\alpha(p),
\]
de modo que $\sigma_p$ mide la curvatura de las secciones normales a $S$ en $p$.
\end{observacion}

\begin{proposicion}
La segunda forma fundamental $\sigma_p$ es \textbf{simétrica} para todo $p\in S$. En consecuencia, $A_p$ y $(dN)_p$ son autoadjuntos respecto a $g_p$, y por tanto \textbf{diagonalizables}.
\end{proposicion}

\begin{notacion}
Los valores propios de $A_p$, ordenados $k_1(p)\leq k_2(p)$, se llaman \textbf{curvaturas principales} de $S$ en $p$; los vectores propios correspondientes $e_1, e_2$ se llaman \textbf{direcciones principales}. Se tiene
\[
K = k_1 k_2, \qquad H = \frac{1}{2}(k_1+k_2),
\]
\[
k_1(p) = \min_{\substack{v\in T_pS\\|v|=1}}\sigma_p(v,v), \qquad k_2(p) = \max_{\substack{v\in T_pS\\|v|=1}}\sigma_p(v,v).
\]
El polinomio característico de $A_p$ es $\lambda^2 - 2H\lambda + K$, y siempre se cumple $K\leq H^2$, con igualdad si y solo si $k_1(p)=k_2(p)$.
\end{notacion}

\begin{definicion}
Un punto $p\in S$ es \textbf{umbilical} (o \textbf{umbílico}) si $k_1(p)=k_2(p)$, equivalentemente $K(p)=H^2(p)$. Si todos los puntos de $S$ son umbilicales, $S$ es \textbf{totalmente umbilical}.
\end{definicion}

\begin{observacion}
Las direcciones de máxima y mínima curvatura son siempre ortogonales entre sí.
\end{observacion}

\begin{observacion}[Curvaturas de superficies elementales]\
\begin{itemize}
\item \textbf{Plano:} $k_1=k_2=0$, $K=H=0$. Totalmente umbilical.
\item \textbf{Esfera $S^2(p_0,R)$:} $A_p = -\dfrac{1}{R}I$, $k_1=k_2=-\dfrac{1}{R}$, $K=\dfrac{1}{R^2}$, $H=-\dfrac{1}{R}$. Totalmente umbilical.
\item \textbf{Cilindro circular $\{x^2+y^2=R^2\}$:} $k_1=-\dfrac{1}{R}$, $k_2=0$, $K=0$, $H=-\dfrac{1}{2R}$. Sin puntos umbilicales.
\end{itemize}
\end{observacion}

\begin{teorema}[Clasificación de superficies totalmente umbilicales]
Sea $S\subset\mathbb{R}^3$ una superficie conexa, orientada y totalmente umbilical. Entonces $S$ es un abierto de un plano (si $k_1=k_2=0$) o un abierto de una esfera $S^2\!\left(p_0,\tfrac{1}{|\lambda|}\right)$ (si $k_1=k_2=\lambda\neq 0$). En particular:
\begin{itemize}
\item Si $S$ es cerrada, es un plano completo o una esfera completa.
\item Si $S$ es compacta, es una esfera.
\end{itemize}
\end{teorema}

\subsection{Curvaturas en coordenadas locales}

\begin{definicion}
Sea $S$ orientada por $N:S\to S^2$. Una parametrización $X:U\to S$ es \textbf{compatible con $N$} si
\[
N\circ X = N^X = \frac{X_u\land X_v}{|X_u\land X_v|}.
\]
\end{definicion}

\noindent
Fijada una parametrización compatible $X$, y tomando $B=\{X_u,X_v\}$ como base de $T_XS$, se definen los \textbf{coeficientes de la primera forma fundamental}:
\[
M = \begin{pmatrix} E & F \\ F & G \end{pmatrix}, \qquad E=|X_u|^2,\quad F=\langle X_u,X_v\rangle,\quad G=|X_v|^2,
\]
y los \textbf{coeficientes de la segunda forma fundamental}:
\[
\Sigma = \begin{pmatrix} e & f \\ f & g \end{pmatrix}, \qquad e = \frac{\det(X_u,X_v,X_{uu})}{|X_u\land X_v|},\quad f = \frac{\det(X_u,X_v,X_{uv})}{|X_u\land X_v|},\quad g = \frac{\det(X_u,X_v,X_{vv})}{|X_u\land X_v|}.
\]
La matriz del endomorfismo de Weingarten respecto de $B$ es $A_X = M^{-1}\Sigma$, y como $\det M = EG-F^2\neq 0$:
\[
A_X = \frac{1}{EG-F^2}\begin{pmatrix} G & -F \\ -F & E \end{pmatrix}\begin{pmatrix} e & f \\ f & g \end{pmatrix}.
\]

\begin{proposicion}[Fórmulas de las curvaturas en coordenadas]
\[
K\circ X = \frac{eg-f^2}{EG-F^2}, \qquad H\circ X = \frac{eG-2fF+gE}{2(EG-F^2)},
\]
\[
k_{1,2}\circ X = H\circ X \mp \sqrt{(H\circ X)^2 - K\circ X}.
\]
Las aplicaciones $K,H:S\to\mathbb{R}$ son diferenciables; $k_1,k_2:S\to\mathbb{R}$ son diferenciables fuera de los puntos umbilicales y continuas en todo $S$.
\end{proposicion}

\begin{proposicion}[Curvaturas de un grafo]
Para $S = \Gr h$ con $h:U\subset\mathbb{R}^2\to\mathbb{R}$ diferenciable y $X(u,v)=(u,v,h(u,v))$:
\[
E = 1+h_u^2,\quad F = h_u h_v,\quad G = 1+h_v^2,\quad EG-F^2 = 1+|\nabla h|^2,
\]
\[
e = \frac{h_{uu}}{\sqrt{1+|\nabla h|^2}},\quad f = \frac{h_{uv}}{\sqrt{1+|\nabla h|^2}},\quad g = \frac{h_{vv}}{\sqrt{1+|\nabla h|^2}},
\]
\[
K\circ X = \frac{\det(\mathrm{Hess}\, h)}{(1+|\nabla h|^2)^2}, \qquad H\circ X = \frac{h_{uu}(1+h_v^2)+h_{vv}(1+h_u^2)-2h_{uv}h_uh_v}{(1+|\nabla h|^2)^{3/2}}.
\]
\end{proposicion}

\begin{observacion}[Invarianza por movimientos rígidos]
Las curvaturas $K$, $H$, $k_1$ y $k_2$ son invariantes por movimientos rígidos de $\mathbb{R}^3$: si $\Phi(x)=Ax+b$ con $AA^t=I$, entonces $K'\circ\Phi = K$, $H'\circ\Phi = H$, etc. (donde las primas indican la curvatura de $\Phi(S)$ con la orientación inducida).
\end{observacion}

\begin{observacion}[Interpretación de $K$ y $H$]\
\begin{itemize}
\item $K>0$ en todos los puntos: la superficie encierra un conjunto convexo (e.g. elipsoides). $K>0$ constante implica que es una esfera.
\item $K=0$: la superficie es localmente un cilindro o un plano.
\item $K<0$: la superficie tiene un punto de silla en cada punto (e.g. el helicoide, con $K<0$ en todos sus puntos).
\item $H=0$: la superficie se llama \textbf{mínima} (busca minimizar área; la curvatura media mide la tensión superficial).
\end{itemize}
\end{observacion}

\begin{proposicion}
Sea $S\subset\mathbb{R}^3$ una superficie compacta. Entonces existe al menos un punto $p\in S$ con $K(p)>0$.
\end{proposicion}

\subsection{Lema y teoremas de Hilbert}

\begin{teorema}[Lema de Hilbert]
Sea $S$ una superficie orientada y $k_1\leq k_2$ las curvaturas principales. Sea $p\in S$ tal que:
\begin{enumerate}
\item $k_1$ tiene un mínimo local en $p$,
\item $k_2$ tiene un máximo local en $p$,
\item $K(p)>0$.
\end{enumerate}
Entonces $p$ es un punto umbilical.
\end{teorema}

\begin{corolario}[Teorema de Hilbert-Liebmann]
Sea $S$ una superficie compacta y conexa con $K$ constante. Entonces $S$ es una esfera.
\end{corolario}

\begin{corolario}[Teorema de Jellett-Liebmann]
Sea $S$ una superficie compacta, conexa, con $H$ constante y $K>0$. Entonces $S$ es una esfera.
\end{corolario}

\end{document}